\documentclass{article}
\usepackage[utf8]{inputenc}
\usepackage{amsmath, amsthm, amssymb, mathpazo, isomath, mathtools}
\usepackage{subcaption,graphicx,pgfplots}
\usepackage{fullpage}
\usepackage{booktabs}
\usepackage{hyperref}
\usepackage{algorithm, algorithmic}
\usepackage{mathtools}
\usepackage{todonotes}

\title{Tarea 1}
%\author{Nicol\'as A Barnafi\thanks{Instituto de Ingeniería Biológica y Médica, Pontificia Universidad Católica de Chile, Chile}, Axel Osses\thanks{Departamento de Ingeniería Matemática, Universidad de Chile, Chile}}
%\author{Nicol\'as A Barnafi}
\date{}

\renewcommand{\vec}{\vectorsym}
\newcommand{\mat}{\matrixsym}
\newcommand{\ten}{\tensorsym}
\newcommand{\der}[2]{\frac{d#1}{d#2}}
\DeclareMathOperator{\grad}{\nabla}
\DeclareMathOperator{\dive}{\text{div}}
\DeclareMathOperator{\curl}{\text{curl}}
\DeclareMathOperator{\tr}{\text{tr}}
\DeclareMathOperator{\sym}{\text{sym}}
\newtheorem{remark}{Remark}
\newtheorem{definition}{Definition}
\newcommand{\R}{\mathbb{R}}
\newcommand{\D}{\mathcal{D}}
\newcommand{\parder}[2]{\frac{\partial\,#1}{\partial\,#2}}

\newcommand{\tin}{\text{in}}
\newcommand{\ton}{\text{on}}

\newtheorem{theorem}{Theorem}
\newtheorem{lemma}{Lemma}
\newcommand{\pts}[1]{[{\bf #1 puntos}] }

\begin{document}

\maketitle

\todo[inline,color=white!90!black]{\textbf{Instrucciones: } La tarea debe ser entregada de manera individual en un informe escrito en \LaTeX a través del buzón habilitado en la plataforma Canvas, donde deben mostrar también el código desarrollado (si es que corresponde). Para su conveniencia, pueden entregar las tareas en un Jupyter Notebook con los desarrollos matemáticos dentro, de modo que sea más cómodo mostrar el código. En cualquier caso, se debe entregar un único archivo como respuesta a la tarea. La política de atrasos será: se calculará un factor lineal que vale 1 a la hora de entrega y 0 48 horas después. Esto multiplicará su puntaje obtenido. Pueden usar ChatGPT u otros modelos solo a conciencia. El uso de salidas de GPT sin su debida comprensión será severamente sancionado.}

\begin{enumerate}
    \item\pts{1} Defina las derivadas de Fréchet de alto orden. Luego, considere el operador $T: L^p(\Omega) \to \R$ dado por $T(f) = \int_\Omega f^p\,dx$ y calcule su derivada de tercer orden.
    \item\pts{2} Enuncie y demuestre la regla de la cadena para derivadas de Fréchet, dada por $d(G\circ F)(u)[h] = dG(F(u))[dF(u)[h]$. Luego, úsela para encontrar la derivada del operador $S:L^\infty(\Omega)\to L^\infty(\Omega)$ dado por $S(u) = \sin u$. Defina de manera precisa todos los espacios usados para generar usar la regla de la cadena. 
    \item\pts{2} El Teorema de Caley-Hamilton establece que toda matriz satisface su propia ecuación característica usando las invariantes: 
                $$ \mat A^3 - I_1(\mat A) \mat A^2 + I_2(\mat A) \mat A - I_3(\mat A) \mat I = 0. $$
                Derive esta ecuación con respecto a $\mat A$ y desarrolle la ecuación para demostrar que $\parder{\det \mat A}{\mat A} = \det (\mat A) \mat A^{-T}$.
    \item\pts{1}  Use la pregunta anterior para demostrar que si $A:\R \to \R^{n\times n}$ es una función matricial, luego
            $$ \der{\det A}{t}(t) = \det A(t) A^{-T}(t) : \der A{t}(t). $$
            \item Si $\phi, \vec u, \ten A$ son un campo escalar, vectorial, y matricial respectivamente en $\R^3$, demuestre usando notación indicial que si $\vec v$ es otro campo vectorial:
                    \begin{itemize}
                        \item\pts{0.5} $$\dive(\phi\vec  u) = \phi \dive \vec u + \vec u \cdot \grad \phi $$
                        \item\pts{0.5} $$\dive(\phi\ten  A) = \phi \dive \ten A + \ten A \grad \phi $$
                        \item\pts{0.5} $$\dive(\vec u \times \vec v) = \vec v\cdot \curl \vec u - \vec u \cdot \curl v $$
                        \item\pts{0.5} $$ \grad(\phi \vec u) = \vec u \otimes \grad \phi + \phi \grad \vec u$$
                    \end{itemize}
    \item Considere el dominio $\Omega=(a,b)$ con $a<b$.
            \begin{itemize}
                \item\pts{1} Considere la función $f:\Omega \to \R$ dada por $f(x) = \sin (2\pi x)$. Para ella, considere un esquema de diferenciación numérica usando el método de diferencias finitas para aproximar su derivada $f'(x)$ en $\Omega$. En particular, calcule la tasa de convergencia teórica para los esquemas de diferencias hacia adelante, hacia atrás, y centradas en la norma $\|\cdot \|_{\infty,h}$. Para ello, considere $a=-1$, $b=1$, y una distancia entre puntos discretos dada por $h  \in \{2^{-n}\}_{n=4}^{20}$, calcule las tasas de convergencia empíricas e interprete los resultados obtenidos.
                \item\pts{1} Repita el ejercicio anterior para la función $f(x) = 2$ e interprete los resultados obtenidos.
            \end{itemize}
        \item\pts{3} Implemente las derivadas discretas de primer orden $\ten D^+, \ten D^-, \ten D$ y construya un ejemplo para validar las tasas de convergencia vistas en clase. Calcule analíticamente la tasa de convergencia de las diferencias centradas para las derivadas de segundo orden $\ten D^-\ten D^+$, y programe un ejemplo para validar también su tasa de convergencia. Finalmente, analice el caso de la composición de diferencias centradas $\ten D\ten D$.
        \item\pts{2} En el apunte se detalla cómo las diferencias centradas se pueden escribir como la composición de los operadores de diferencias hacia adelante y hacia atrás según $\ten D^+ \ten D^- = \ten D^-\ten D^+$. Demuestre teóricamente que en cambio $\ten D^+\ten D^+$ (o $\ten D^-\ten D^-$) no es buena aproximacion de la derivada de segundo orden.
    \item Considere el dominio $\Omega=(0,1)$. En él, discretizar la ecuación 
            $$ \mathcal L u = 1 $$ 
            con el operador diferencial $\mathcal L$ dado por
            $$ \mathcal L u = - u'' + bu' + cu, $$
            con condiciones de borde dadas por $u(0) = 1$, $u(1) = 2$. Este problema se conoce como la ecuación de Advección-Difusión-Reacción (ADR).  Calcularemos la sensibilidad del problema a sus parámetros en términos del condicionamiento del sistema, que está intrínsicamente relacionado con la dificultad numérica del problema. Llamaremos a la matriz inducida por la discretización $\mat L_h$. 
            \begin{itemize}
                \item\pts{1} Muestre la ecuación correspondiente al nodo $i$ de la malla para (i) diferencias hacia adelante para $u'$ y (ii) diferencias centradas para $u'$. 
                \item\pts{2} En cada una de las dos discretizaciones del punto anterior, muestre que la discretización obtenida es consistente y calcule el orden de consistencia.
                \item\pts{1} Grafique en un gráfico log-log el condicionamiento de la matriz $\mat L_h$ con respecto a los siguientes valores: $h \in \{2^{-n}\}_{n=2}^8$, $b \in \{0.01,0.1,1,10,100\}$, y $c \in \{0.01,0.1,1,10,100\}$. Hágalo solo para el esquema donde $u'$ es aproximado con diferencias centradas.
                \item\pts{2} Demuestre teóricamente usando el Teorema de Equivalencia de Lax que la solución discreta obtenida con la matriz $\mat L_h$ converge a la solución continua si $b=0$. Para ello, extienda la demostración del apunte para el Laplaciano.
                \item\pts{2} Usando el método de las soluciones manufacturadas, muestre que la tasa de convergencia numérica a una solución analítica es la esperada por la teoría. 
            \end{itemize}

\end{enumerate}

\todo[inline,color=white!90!black]{\textbf{Nota: } Abriremos un foro en Canvas para revisar cualquier typo y/o error que haya en el enunciado.}
\end{document}

